\chapter{基于統計學習的行人檢測框架}

\section{特征提取}
注釋可採用腳注或章節附注的方式，並按照本學科通行範式逐一注明需解釋的詞句、資料來源或版權許可等。

\subsection{HOG 特征}
每個表均應包含表序號與表題，表題置於表上方；表應在正文中先提及後出現。\cref{tab:video-info} 給出了一個表格示例。

\begin{table}[htbp]
  \centering
  \caption{檢測視頻的效果信息}
  \label{tab:video-info}
  \vspace{0.2cm}
  \begin{tabular}{lccc}
    \toprule
    視頻序列 & 每幀大小 & 平均檢測用時 & 平均正確率 \\
    \midrule
    Video 1 & 640 x 480 & 62 毫秒/幀 & 93.3\% \\
    Video 2 & 640 x 480 & 65 毫秒/幀 & 92.7\% \\
    Video 3 & 320 x 240 & 44 毫秒/幀 & 94.1\% \\
    \bottomrule
  \end{tabular}
\end{table}
